\documentclass[12pt]{article}
\usepackage[utf8]{inputenc}
\usepackage[T1]{fontenc}
\usepackage{hyperref}
\title{Wapichan Language Resources for K-12 Curriculum Development}
\author{Compiled by Claude Sonnet 5 (claude-sonnet-5) for Andr\'as Kornai}
\date{}
\begin{document}
\maketitle
\section{Introduction}
Wapichan (ISO 639-3 wap; also Wapishana, Wapixana, Wapichana) is an Arawakan language of Guyana and Brazil with about 13,000 speakers. Searching under all spellings is advisable.
\section{Primary Educational Resources}
\begin{itemize}
\item \emph{Gram\'atica Pedag\'ogica da L\'ingua Wapichana} (Museu do Índio/FUNAI, 2023, 302 pp., Brazilian variety, made with Wapichana teachers): \url{https://www.gov.br/museudoindio/pt-br/centrais-de-conteudo/livros-e-publicacoes/materiais-didaticos/arquivos/GramaticaWapichana.pdf}
\end{itemize}
\section{Linguistic and Reference Materials}
\begin{itemize}
\item Wapishana language, overview: \url{https://en.wikipedia.org/wiki/Wapishana_language}
\item ``On Making A Wapichan Dictionary'' (essay by Pauline Melville): \url{https://3quarksdaily.com/?p=36529}
\item Demonstratives, determiners and definiteness in Wapichana (article): \url{https://revistas.ufpr.br/letras/article/view/50821}
\item Wapishana numerals: \url{https://omniglot.com/language/numbers/wapishana.htm}
\end{itemize}
\section{Religious and Community Resources}
Wikipedia lists the Guyanese organization Wapichan Wadauniinao Ati'o for language preservation; I found no website for it. Povos Ind\'igenas no Brasil: \url{https://pib.socioambiental.org/en/povo/Wapichana}
\section{Digital and Audio Resources}
None found in search results.
\section{Curriculum Development Context}
Wapichan is not known to be taught in schools. The Brazilian pedagogical grammar is the most curriculum-oriented item found; no freely available Guyanese dictionary was located.
\end{document}